\documentclass[11pt]{article}
\usepackage[T1]{fontenc}
\usepackage{amsmath,amssymb,amsthm}
\usepackage{enumitem}
\usepackage[margin=1in]{geometry}
\usepackage{hyperref}

\newtheorem{theorem}{Theorem}

\newcommand{\XiBC}{\Xi^{BC}_{l,a}}
\newcommand{\Aeta}{A_\eta}
\newcommand{\Keta}{K_\eta}
\newcommand{\Rhofin}{\mathcal R_{\mathrm{fin}}}

\begin{document}

\title{Step 166: Cascade Diagnostic-Complete RH Membrane Packet Summary}
\author{}
\date{}
\maketitle

\section{Cascade Typed Context}

The orientation is adequacy. The inherited parent residual is
\[
\XiBC=B^*_{l,a}\Pi_{Y_a}B_{l,a},
\]
with machine label \(\mathrm{Xi\_BC}=B^*\Pi_YB\). The global active operator is
\[
C_lP_\eta .
\]
The carrier is the Burnol/Sonine pulled-evaluator carrier. The cascade has three sibling child branches under \(\mathrm{Xi\_BC}\).

\subsection{Branch A: Calkin-Bridge Lane G2--G5}

The typed Calkin context is
\[
\Aeta=C^*(P_\infty,M_{m_l},P_\eta,I),\qquad
\Keta=\Aeta\cap K(H),
\]
with quotient context \(\Aeta/\Keta\). The carried gates are G1 framework-defined, G2 faithful symbol open external proof, G3 normal form \(C_l=U_lW_l+K_l\) open external proof, G4 lower-faithfulness of \(U_l\) open external proof, and G5 compact remainder \(K_l\in K(H)\) open external proof. The branch status is \(\mathrm{split\_external\_theorem}\), sourced from step 163.

\subsection{Branch B: Per-Zero Finite-Carrier Lane SL164.1}

The typed context is the per-zero stratification
\[
H_\eta=\widehat{\bigoplus}_{\rho}H_{\eta,\rho},\qquad
A_{\eta,\rho}\subseteq A_\eta .
\]
The finite carrier is
\[
\Rhofin=\{\hbox{first three zeros}\},\qquad
A_{\mathrm{fin}}=\{1/2\},\qquad
K_{\mathrm{fin}}=\{0\}.
\]
The active sub-residual is
\[
\Xi_{\mathrm{matrix\_source},l,\Rhofin,A_{\mathrm{fin}},K_{\mathrm{fin}}},
\]
with parent \(\mathrm{Xi\_BC}\). The dormant residuals are
\(\Xi_{\mathrm{coupling},l,\rho,\rho'}\), activated by off-diagonal coupling computation;
\(\Xi_{\mathrm{summability},l}\), activated by diagonal structure computation; and
\(\Xi_{\mathrm{orth},\eta,\rho,\rho'}\), activated if SL164.2 is rejected. SL164.1 has status support-only, SL164.2 has status working hypothesis, and the branch status is \(\mathrm{finite\_carrier\_diagnostic\_indeterminate}\), sourced from step 164.

\subsection{Branch C: Shifted Co-Poisson Exact-Route Lane H1--H5}

The typed context is the shifted Burnol packet
\[
u_{l,a}=B_{l,a}u
\]
with candidate factorization
\[
M(Bu_{l,a})(s)=\zeta(s)\alpha_{l,a,u}(s).
\]
The child residual is \(\Xi_{\mathrm{cP\_shifted},l}\), with parent \(\mathrm{Xi\_BC}\). The carried gates are H1 strip/packet context framework-defined, H2 factorization existence open external proof, H3 typed properties of \(\alpha_{l,a,u}\) open external proof, H4 endpoint/zero admissibility open external proof, and H5 exact implication to \(\mathrm{Xi\_BC}\) closure open external proof. The branch status is \(\mathrm{split\_external\_theorem}\), sourced from step 165.

\subsection{Retained No-Gos}

The retained framework no-gos are public-shadow non-promotion and finite-window Calkin blindness. They foreclose public-shadow promotion routes and finite-window or finite-section completed-Calkin certificates.

\section{Cascade Summary Theorem}

\begin{theorem}[RH membrane packet cascade diagnostic-complete on the Burnol/Sonine pulled-evaluator carrier]
On the inherited Burnol/Sonine pulled-evaluator carrier with active residual
\(\XiBC=B^*_{l,a}\Pi_{Y_a}B_{l,a}\) and active operator \(C_lP_\eta\), the framework's analytical structural cascade is diagnostic-complete at three child branches:
\begin{enumerate}[label=(\alph*)]
\item Branch A: the Calkin-bridge route requires the split bridge theorem G2--G5 in the typed Calkin context \(\Aeta/\Keta\).
\item Branch B: the per-zero finite-carrier route requires external Burnol/Sonine projected-kernel evaluator records SL164.1, with active source residual \(\Xi_{\mathrm{matrix\_source},l,\Rhofin,A_{\mathrm{fin}},K_{\mathrm{fin}}}\).
\item Branch C: the shifted co-Poisson exact-route requires the split factorization theorem H1--H5 with \(M(Bu_{l,a})(s)=\zeta(s)\alpha_{l,a,u}(s)\) and its exact implication.
\item The two retained framework no-gos foreclose public-shadow promotion routes and finite-window or finite-section completed-Calkin certificates.
\item Under the inherited records through step 165, no further refinement within any single branch on this carrier closes \(\mathrm{Xi\_BC}\); closure requires either external content satisfying one of A, B, or C, or a pivot to a different primary carrier such as the Hecke carrier, character-source frame, layer-dissolving framework, or duality-confinement framework.
\end{enumerate}
\end{theorem}

\begin{proof}
The proof is compositional. Step 163 supplies Branch A as the Calkin-bridge split external theorem in the typed Calkin context. Step 164 supplies Branch B as the per-zero finite-carrier diagnostic record, including SL164.1 support-only status, SL164.2 working-hypothesis status, the matrix-source residual, and the dormant residual activation rules. Step 165 supplies Branch C as the shifted co-Poisson split external theorem with candidate zeta factorization and exact-route gates. Steps 158--162 carry the public-shadow non-promotion and finite-window Calkin blindness no-gos. The diagnostic-complete versus proof-complete distinction follows the inherited cascade rule from cheat sheet section 9: the framework output here is the typed cascade structure, while the proofs of G2--G5, SL164.1, and H1--H5 remain external.
\end{proof}

\section{Typed External-Content Interface}

\subsection{Branch A Interface}

External work must respect \(\Aeta=C^*(P_\infty,M_{m_l},P_\eta,I)\), \(\Keta=\Aeta\cap K(H)\), and the quotient \(\Aeta/\Keta\). The obligations are G2 faithful symbol, G3 normal form \(C_l=U_lW_l+K_l\), G4 lower-faithfulness of \(U_l\), and G5 \(K_l\in K(H)\). A successful external theorem would make the Branch A route sufficient for its targeted \(\mathrm{Xi\_BC}\) closure implication. It must not use public-shadow promotion or finite-window Calkin evidence as a completed certificate.

\subsection{Branch B Interface}

External work must respect \(H_\eta=\widehat{\oplus}_{\rho}H_{\eta,\rho}\), \(A_{\eta,\rho}\subseteq A_\eta\), \(\Rhofin\), \(A_{\mathrm{fin}}=\{1/2\}\), and \(K_{\mathrm{fin}}=\{0\}\). The obligations are SL164.1 external projected Burnol/Sonine kernel evaluator records and evaluation of \(\Xi_{\mathrm{matrix\_source},l,\Rhofin,A_{\mathrm{fin}},K_{\mathrm{fin}}}\), while preserving dormant activation rules. A successful external record would decide whether this route supplies its targeted \(\mathrm{Xi\_BC}\) closure implication or identifies the next activated residual. It must not promote support-only SL164.1 data or substitute finite-window Calkin evidence for a completed theorem.

\subsection{Branch C Interface}

External work must respect \(u_{l,a}=B_{l,a}u\), \(M(Bu_{l,a})(s)=\zeta(s)\alpha_{l,a,u}(s)\), and \(\Xi_{\mathrm{cP\_shifted},l}\) under parent \(\mathrm{Xi\_BC}\). The obligations are H1 strip/packet context, H2 factorization existence, H3 typed properties of \(\alpha_{l,a,u}\), H4 endpoint/zero admissibility, and H5 exact implication. A successful external theorem would make the Branch C route sufficient for its targeted \(\mathrm{Xi\_BC}\) closure implication. It must not bypass endpoint or zero admissibility gates, and it must not reclassify public-shadow or finite-window evidence as theorem-grade content.

\section{Cascade Diagnostic-Complete Declaration}

Step 166 declares diagnostic-complete status at the cascade level. This differs from step 163, which declared diagnostic-complete status at the current carrier level for a single Calkin-bridge branch. Step 166 records the three-branch cascade under the common parent residual \(\mathrm{Xi\_BC}\).

\section{Nonclaim Boundary}

This step does not prove RH. It does not prove \(\mathrm{Xi\_BC}=0\). It does not close any of A, B, C. It does not close G2--G5, SL164.1, or H1--H5. It does not remove either retained no-go. It does not pivot to a different carrier. It does not begin step 167 lanes. It does not supersede or merge the three branches. It does not promote SL164.1 or SL164.2 to accepted bridge source. It does not make \(\mathrm{Xi\_BC}\) closure conditional on a single external lane: any one of A, B, or C closing would suffice for its carried route, but none is established here.

\end{document}
